\section{Discusión}

Los resultados obtenidos para la densidad experimental del fluido,
$\rho_f = (0{,}991 \pm 0{,}078)\,\text{g/cm}^3$, con un error porcentual
de $0{,}86\%$ y una incertidumbre relativa de $7{,}89\%$,
indican que el intervalo experimental
$0{,}913 \leq \rho_f \leq 1{,}069\,\text{g/cm}^3$ contiene al valor de
referencia del agua ($1{,}00\,\text{g/cm}^3$). Esto implica que, dentro
de los márgenes de incertidumbre del montaje, el modelo del principio
de Arquímedes describió adecuadamente la relación entre el empuje y el
volumen de fluido desplazado, cumpliéndose así los objetivos de medir
el empuje sobre la varilla y verificar experimentalmente dicho
principio.

Sin embargo, el ajuste lineal arrojó un intercepto
$b = (0{,}227 \pm 0{,}068)\,\text{g}$, cuyo valor no resultó
despreciable frente a su incertidumbre ($b/\Delta b \approx 3{,}3$), a
pesar de haberse asumido $b \approx 0$ en el análisis. Este
desplazamiento sistemático pudo originarse en un ligero descalibre del
punto de referencia de la balanza o en efectos de tensión superficial
en el menisco al momento del primer contacto de la varilla con el
agua. Las mediciones realizadas no permiten atribuir el desplazamiento
a una fuente única, por lo que estas posibilidades deben verificarse
en una repetición del montaje.

En la Parte 2, el error porcentual promedio entre los tiempos de
vaciado experimentales y teóricos fue de $14{,}65\%$, y en cuatro de
las cinco alturas medidas ($h = 10$, $8$, $6$ y $4\,\text{cm}$) el
tiempo experimental resultó sistemáticamente mayor que el tiempo
teórico $t_v$. Este patrón, más que dispersión aleatoria, es
indicativo de un error sistemático consistente con las limitaciones
del modelo de Torricelli, el cual asume flujo ideal sin fricción; en
la práctica, pérdidas de energía por fricción y efectos viscosos en el
orificio real reducen la velocidad de salida por debajo del valor
ideal $\sqrt{2gh}$, prolongando así el tiempo real de vaciado respecto
al teórico. El único caso con signo invertido, $h = 2\,\text{cm}$
(error de $5{,}73\%$), sugiere que a alturas bajas el error de
reacción humana al operar el cronómetro puede adquirir un peso
proporcionalmente mayor sobre tiempos de vaciado más cortos, pudiendo
enmascarar parcialmente el sesgo sistemático. En la misma línea, el ajuste potencial
$t = (88{,}72)h^{0{,}651}$ mostró un exponente experimental
$n = 0{,}651$, superior al valor teórico $n = 0{,}5$ predicho por
Torricelli, lo que confirma que la relación entre $t$ y $h$ es
sublineal, como anticipa la teoría, aunque el modelo ideal no
reproduce con exactitud la magnitud de dicha dependencia.

En conjunto, estos resultados permiten verificar cualitativamente el
principio de continuidad, dado que la tendencia esperada entre el
tiempo de vaciado y la altura se observó consistentemente
($R^2 = 0{,}9806$). No obstante, la verificación cuantitativa resultó
menos precisa que en la Parte 1, reflejando que las aproximaciones de
flujo ideal del teorema de Torricelli tienen un alcance limitado
frente a las condiciones reales del experimento.
